\section{Ampliação ARTESP e conferência territorial de Mauá em 2 de novembro}
O pacote original \texttt{artesp\_gtfs.zip}, obtido no portal oficial da ARTESP \citep{artespgtfs2026}, registra versão \texttt{20261007} em \texttt{feed\_info.txt}. Seus calendários WD, SAT e SUN abrangem 1 de janeiro a 31 de dezembro de 2026, sem exceções de feriados. O arquivo preservado possui 46.662.768 bytes, MD5 \nolinkurl{8c26295e6895e38b407fd516520ccf4f} e SHA-256 \nolinkurl{f3fdca9599e512e191f1bb78cebcafab1bac8a458fbdd5c5e71cc2d10d97d45a}. A URL específica do recurso, seus metadados e a conferência dos CSV originais constam de \nolinkurl{auditoria-artesp-maua-2026-11-02/original-gtfs-verification.json}. O hash identifica o arquivo efetivamente examinado, independentemente de mudanças posteriores no portal.

O comando \texttt{import\_artesp} combina os registros com a base SPTrans sem substituir o catálogo histórico simulado de Mauá. O prefixo \texttt{artesp:} distingue identificadores das novas entidades e evita colisões com os existentes. A base combinada \texttt{network.sqlite3} possui SHA-256 \nolinkurl{f772ebf8b55b1397d1dadc8624449eab33b4080625d57b51b7d0810af3f1317f}. Sua consulta é somente leitura. A materialização compacta combinada concluiu as etapas de custo 0 e 1; a configuração \texttt{SPTRANS\_LEGACY\_DATABASE} mantém o artefato anterior como fonte complementar, sem atribuir suas etapas mais profundas aos novos pares ARTESP. As buscas publicam alternativas progressivamente. Os registros ARTESP não são enviados ao Olho Vivo como se fossem linhas SPTrans: ausência de previsão ao vivo deve permanecer explícita.

O recorte municipal usa a malha detalhada IBGE 2024, arquivo \texttt{SP\_Municipios\_2024.zip}, SHA-256 \nolinkurl{92a54d06e6f5a9749368e522861b785eae8c797329c9fa09cbb1a0696319d8e3}, no sistema SIRGAS 2000 (EPSG:4674). A geometria de Mauá, código \texttt{3529401}, foi preservada em GeoJSON. O predicado espacial inclui a borda; a distância à divisa foi medida na projeção UTM 23S (EPSG:31983). Não foram fundidas paradas próximas ou de sentidos opostos. Contaram-se os identificadores de viagens uma vez por data de serviço, após verificar toda a sequência temporal, e não uma vez por passagem em cada parada. Os 1.491 registros selecionados passaram nessa conferência. Dias úteis não representam automaticamente feriados que ocorram de segunda a sexta-feira.

As consultas de referência usam os pares \nolinkurl{artesp:48804} a \nolinkurl{artesp:5776}; \nolinkurl{artesp:64476} a \texttt{18859}; \nolinkurl{artesp:7700} a \texttt{1010053}; e \texttt{18935} a \texttt{18859}. Foram executadas pelo módulo de produção \nolinkurl{staged_adapter.progressive}, com custo máximo 1, sem substituir serviços ou horários por respostas simuladas. O conjunto de evidências contém os parâmetros, instantes de execução, horários consultados, respostas integrais, identificadores de viagens, durações, caminhadas e hashes dos módulos. A assinatura dos arquivos de banco, por tamanho e instante de modificação, permaneceu inalterada durante as consultas. Os tempos locais não incluem a comunicação até o celular, e a amostra não representa todos os pares metropolitanos.

Para reprodução, o diretório \nolinkurl{auditoria-artesp-maua-2026-11-02/} reúne o censo por parada, linha e viagem, a relação de origens materializadas, os scripts de conferência e consulta, a revisão Git e as alterações locais pertinentes. A base acadêmica anterior e os PDFs precedentes foram preservados. A demonstração funcional não substitui conferência de operação, acessibilidade, atualização periódica dos dados ou avaliação com passageiros. A disponibilidade ferroviária vem do GTFS SPTrans preservado, com frequências e estimativas próprias; suas viagens não foram somadas aos totais de ônibus ARTESP.
